% PASS20260906 -- spread-ladder reversible GC optimality and scheduler binding
\section*{Current VM frontier: checkpoint placement is geometry, not policy}
\addcontentsline{toc}{section}{Current VM frontier: checkpoint placement is geometry, not policy}

\noindent\textbf{Scope.}
This insert records an exact finite-geometry optimality theorem for the reversible
virtual machine and its binding to the adaptive scheduler.  It is a software/control
result: hop-crossing counts on the certified $W(3,3)$ substrate, not Joules, optical
loss, or a fabricated device.

\paragraph{The theorem.}
The reversible runtime executes compute--copy--uncompute with STRONG Merkle
checkpoints.  Releasing a checkpoint tier is a boundary-crossing operation, so
checkpoint \emph{placement} has a geometric floor.  Exact integer certification of
\[
A^2 = 8I - 2A + 4J
\]
for the $W(3,3)$ adjacency matrix gives the Laplacian spectrum $0^1,10^{24},16^{15}$,
hence for every $m$-vertex subset $S$,
\[
|\partial S| \;\ge\; \frac{m(40-m)}{4}.
\]
All $36$ spreads were re-derived by exact cover, and for every ordered spread the ten
nested rungs of $4i$ points attain this bound with equality at every tier
($i=1,\dots,10$; all $1023$ nonempty line unions per spread verified).  Therefore the
spread ladder is the boundary-optimal reversible checkpoint hierarchy:
\[
\boxed{
\text{per-tier release } 4i(10-i),\qquad
\text{peak live boundary } 100 \text{ at the } 20|20 \text{ half-ladder},\qquad
\text{full cycle } 2\cdot 660 = 1320 \text{ crossings},
}
\]
with zero retained-point migration on resize.  The peak value is exactly the certified
minimum boundary of the $20|20$ cut, so a half-checkpointed reversible guest sits on
the substrate's spectral Cheeger value; no placement can do better at half retention.

\paragraph{The binding.}
The adaptive scheduler chooses \emph{how much} to retain (segment size $B$, recursion
levels $L$) from live memory pressure; the ladder now determines \emph{where}.  A
chosen strategy pins $L{+}1$ STRONG checkpoint roots on ladder rungs in retention
order at release cost $4i(10-i)$ each; strategies needing more than ten live roots
degrade the overflow to HASH\_ONLY audit strength---content identity without byte
retention, the strength distinction the temporal GC already defines.  No root is
silently upgraded or discarded, and the destructive-discard control remains excluded
from the reversible frontier.

\paragraph{Artifacts and verification boundary.}
The optimality certificate is
\texttt{analysis/w33\_spread\_ladder\_reversible\_gc.py} with frozen artifact
\texttt{data/w33\_spread\_ladder\_reversible\_gc.json}
(\texttt{sha256:ebb04e0ed67e3c014081d19ce5fe867f3e29379e7912a159130cbdb20d0d6117});
the scheduler binding is
\texttt{analysis/w33\_ladder\_checkpoint\_placement.py} with
\texttt{data/w33\_ladder\_checkpoint\_placement.json}
(\texttt{sha256:0cf744ccc6fb893c062cc66ae9878f83228cc2d6eb561659f46b924bac5e32fc}).
The Holotrade control-plane object is
\texttt{scheduler/w33-reversible-checkpoint-ladder.js}, fail-closed
(\texttt{dispatchable=false}) pending runtime topology attestation.  The result closes
the placement half of reversible scheduling; it does not supply measured energy,
physical routing, or a hardware claim.

\paragraph{Corollary: hypervisor suspension is rung 9.}
When the $1296$-state fibre hypervisor ($216\times_{36}216=36\times6\times6$) suspends
one $216$-state guest fork, the state that must be retained is exactly the shared
$36$-state base coordinate --- and a $36$-point region is rung $9$ of the ladder, with
release boundary $36=36\cdot4/4$, the spectral minimum.  Suspend/resume therefore costs
$72$ boundary crossings, the minimum possible for any $36$-point retention on $W(3,3)$,
and nesting inside rung $10$ gives zero retained-point migration.  The base is shared,
so hypervisor control state adds no second region.  Certificate:
\texttt{analysis/w33\_fibre\_suspension\_ladder.py},
\texttt{data/w33\_fibre\_suspension\_ladder.json}
(\texttt{sha256:41fd5c87e42c2b759d37727b1c062cb59f28ac8f3b720aec998a7e4b5721426e}).
